\begin{reviewcomment}{6}
	Provide additional justification for the PASS/FAIL leakage criteria.
	In the manuscript, a failure is defined only if the mean AUC is greater than 0.55 and at least 50\% of the 30 seeded permutation tests are significant.
	This leads to a criterion based on reproducibility, although the exact thresholds seem to be application-defined.
	The authors need to better justify the choice of 0.55 and 50\% and ideally perform a sensitivity analysis to show that the overall conclusions are robust to reasonable alternative thresholds.
\end{reviewcomment}

\begin{reviewresponse}{6}
	The methodology now justifies both thresholds.
	The AUC threshold of $0.55$ is an effect-size floor that keeps a significant but negligible AUC excess from triggering a FAIL.
	The $50\%$ threshold requires significance in at least $15$ of the $30$ policy-disjoint splits, far above the $5\%$ rejection rate expected per split under the null, so a FAIL reflects an effect that recurs across splits.
	In the evaluated run, the stability condition is the binding one: the AUC condition alone would flag three additional PCH and three additional SPP tasks, each significant in fewer than half of the splits.
	A new sensitivity analysis re-applies the criterion over AUC thresholds of $0.52$--$0.60$ and significant-seed fractions of $0.3$--$0.7$.
	At fractions of at least $0.5$, at most the same two tasks fail; lower fractions add at most two SPP contrasts, and the PI verdict never changes.
	A null calibration with policy-independent gas features gives a spurious-failure rate of about $3\%$ per PCH task and none for SPP.
	The overall conclusions are therefore robust to reasonable alternative thresholds.

	\RevisionLocation{Section~6.2.1, paragraph ``Operational thresholds and sensitivity'' and the two paragraphs that follow (\Loc{ev-thr}).}
\end{reviewresponse}
